% Response to reviewers for the revised PoolParty manuscript.
\documentclass[11pt]{article}

\usepackage[margin=1in]{geometry}
\usepackage[T1]{fontenc}
\usepackage[utf8]{inputenc}
\usepackage{lmodern}
\usepackage{xcolor}
\usepackage{enumitem}
\usepackage{parskip}
\usepackage[round]{natbib}
\usepackage{xr-hyper}
\usepackage[hidelinks]{hyperref}
\externaldocument[ms-]{manuscript_clean}

\definecolor{responseblue}{HTML}{0070C0}
\definecolor{noticegray}{HTML}{555555}
\newenvironment{response}{\par\color{responseblue}}{\par}
\newenvironment{msquote}
  {\begin{quote}\color{responseblue}\itshape\noindent\ignorespaces}
  {\end{quote}}
\providecommand{\bibcommenthead}{}
\setlist[enumerate,1]{leftmargin=*,itemsep=1.2\baselineskip,topsep=\baselineskip}
\newcommand{\reviewersection}[1]{%
  \bigskip\par\noindent{\large\bfseries #1}\par\medskip\hrule\medskip}

\begin{document}

\begin{center}
{\Large\bfseries Response to Reviewers}\\[0.6em]
{\large PoolParty: streamlined design of DNA sequence libraries in Python}\\[0.8em]
Zhihan Liu, Aidan Cordero, Justin B. Kinney\\[0.7em]
\end{center}

\begin{response}
We thank the Editor and Reviewers for their thoughtful critiques, to which we provide a point-by-point response below. The manuscript has been improved substantially as a result of this feedback. Additional edits have also been made throughout the manuscript to improve clarity. 
\end{response}

\reviewersection{Editor Comments}

In general, the reviewers had positive impressions of PoolParty, and it can
likely be revised to a state that will be acceptable. Please carefully address
each point raised by the reviewers. It's especially important to provide
benchmarking data, comparisons to existing tools, and not to overstate the
benefits the tool may provide or fail to discuss limitations. Reviewer 1
recommends comparing PoolParty to 2--3 other comparable tools. However, rather
than arbitrarily picking two or three other tools for comparisons, you should
use a well-defined process for picking which tools, and how many, to include in
comparisons. Perhaps 2--3 is adequate; however, that may be too few tools for
comparisons. The comparisons should include, at a minimum, benchmarking (e.g.,
runtime, memory usage), program features, and similarities/differences in
program outputs.

\begin{response}
We now provide a more thorough comparison between PoolParty and other available tools for desinging sequence libraries, including in the Introduction, Discussion and a new Supplemental Information Supplemental document. Tables S2 and S3 compare PoolParty to a wide range of available tools, and Figs.~S3 and S4 compare PoolParty to its two closest competitors VaLiAnT and tangermeme, in specific use cases. We additionally attempted to provide a use case comparison to MPRAnator, but this tool runs through a web interface that is no longer functional. We have also benchmarked the compute and memory requirements of PoolParty on the three examples feature in the main text (Fig.~S2 and Table S1). Resource usage is shown only for PoolParty, as none of the competing tools were able to generate the specific libraries shown in those examples. 
\end{response}

\reviewersection{Reviewer 1}

\begin{enumerate}
\item In the GB1 example, the library exceeds 540,000 sequences, yet the paper
provides no runtime or memory consumption benchmarks. Readers cannot assess the
tool's practical usability across different library scales. The authors should
supplement benchmarking data, including runtime and peak memory usage across
libraries of increasing sizes (e.g., 1K, 10K, 100K, 500K sequences), and specify
the test environment configuration.

\begin{response}
The revised text benchmarks the resourcee requirements of PoolParty for each of the three examples featured in the main text (in Figs.~1, 2, and 3), using libraries of increasing size. The results are reported in the new Fig.~S1 and Table S1. The new Supplemental Information file additionally reports the machine configuration uses in these benchmarking tests. This new material is referenced in the Discussion. 
\end{response}

\item The paper highlights design cards as a core feature, emphasizing their use
as covariates in downstream analyses. However, it does not evaluate whether
design card variables provide significant incremental predictive value beyond
sequence-derived features. The authors should either supplement a comparative
experiment quantifying the added information gain of design cards, or explicitly
state in the Discussion that this validation is planned as future work.

\begin{response}
Design cards record the specific design choices used to construct each sequence. We claim in the paper that this can facilitate surrogate modling studies by making these design choices readily available as covariates, and support this claim by illustrating in Fig.~3 an sample analysis that does this. The does not, however, claim that the information in design cards increases the predictive performance of sequence-to-function models. 
\end{response}

\item The introduction mentions VaLiAnT and MPRAnator, noting their limitations,
but does not provide a systematic comparison between PoolParty and these tools.
The authors should provide a table comparing PoolParty with 2--3 existing tools
on core features, and demonstrate code conciseness in at least one common use
case, to help readers objectively assess PoolParty's advantages.

\begin{response}
Tables S2 and S3 in the new Supplemental Information document provide a systematic comparison between PoolParty and other competing tools. Figs.~S3 and S4 further provide code comparisons between PoolParty and the two competing packages, VaLiAnT and tangermeme, in two specific use caes. As mentioned above, we additionally tried to provide a comparison to MPRAnator, but that software (which runs through a web interface) is no longer functional. The Introductiona nd Discussion have also been revised to more explicitly compare PoolParty to existing tools. 
\end{response}

\item Provide benchmarking data on runtime and memory usage across libraries of
varying sizes.

\begin{response}
This is addressed in the response to Comment 1 above. 
\end{response}

\item Include a table comparing features with 2--3 existing tools.

\begin{response}
This is addressed in the response to Comment 3 above. 
\end{response}

\item To further contextualize this research within the existing academic
landscape and strengthen its theoretical support, it is recommended that the
authors supplement several relevant representative references.

[1] \url{https://doi.org/10.1093/nsr/nwaf495}

[2] \url{https://doi.org/10.1038/s41467-026-72882-y}

[3] \url{https://doi.org/10.1186/s13059-025-03606-6}

[4] \url{https://doi.org/10.1016/j.patcog.2025.112078}

\begin{response}
Having read these papers, we do not see how these references are relevant to PoolParty for the more general problem of sequence library design: the first paper discusses multimodal representations of small molecules for drug discovery; the second paper models single-cell transcriptomic and surface-protein measurements; the third paper discusses read mapping to pangenome graphs; and the fourth paper concerns circRNA-miRNA interaction prediction. We therefore decided to to cite these papers in the manuscript. 
\end{response}

\item Either provide empirical analysis of design cards' incremental
contribution to surrogate modeling performance, or explicitly discuss this as
future work.

\begin{response}
This is addressed in the response to Comment 2 above. 
\end{response}
\end{enumerate}

\reviewersection{Reviewer 2}

\begin{enumerate}
\item Additional statistical readouts describing pool generation would be
useful. For example, how many unique sequences vs duplicates were produced in
each pool (i.e. out of the universe of all permutations that meet the
specifications)? How unique are the sequences (min/max/average hamming distance)?
How frequent are homopolymers? Etc.

\begin{response}
This is a good suggestion. To provide this information have added a \texttt{.stats()} method to the Pool class. This method computes a number of summary statistics about the pool, including the size of the design spzce (when this is fixed), the number of unique/duplicated sequences generated, the pairwise Hamming distances between generated sequences, the frequencies of homolopymer occurances, etc. We have added a short paragraph to the Sequence metadata subsection describing this method, as well as documentation at  at \url{https://poolparty.readthedocs.io/en/latest/pool.html#summarising-a-library-stats}. 
\end{response}

% \begin{response}
% We agree that these statistics are useful for checking the generated library.
% We have added a \texttt{.stats()} method that can be called on any DNA Pool in
% the DAG, summarizing the size of the specified design space when it is fixed,
% the number of unique/duplicated sequences generated, pairwise Hamming
% distances, homopolymer frequencies, and several other basic properties of the
% Pool. We have also added documentation and examples at
% \url{https://poolparty.readthedocs.io/en/latest/pool.html\#pool-stats}. We have
% added a short paragraph to the Sequence metadata subsection of the
% Implementation section:
% \end{response}

% \begin{msquote}
% PoolParty also provides a \texttt{stats} method for summarizing a DNA Pool,
% including its unique and duplicate sequence counts, pairwise Hamming distances,
% and homopolymer frequencies.
% \end{msquote}

\item The authors disclose that LLMs were used to develop, debug, and document the code. Did the authors also write a test harness, and what edge cases were nominated by the LLMs for testing vs what edge cases were nominated by the human developers for testing?

\begin{response}
Yes, the PoolParty codebase includes over 3{,}000 tests that can be run using pytest. In our experience, the coding agents were quite good at testing mechanical aspects of the code (e.g., invalid inputs and boundary conditions). However, human-written use cases (from which most of the manual tests were derived) were essential for testing the types of inputs and requests likely to be encountered in real biological applications. The current suite consists of a mixture of these two types of tests. We did not track the origin of each test, however, and so precisely quantify how many came from each source.  
\end{response}

% \begin{response}
% Yes. PoolParty has over 3,000 tests written with pytest. These tests are from a
% combination of humans and coding agents. We did not track the origin of every
% individual test, so we cannot separate them exactly after the fact. In our
% experience, coding agents are particularly good at testing mechanical aspects
% of the code (e.g., invalid inputs and boundary conditions) instead of
% biologically realistic tests (e.g., DMS and MPRA designs that respect
% biological and physical constraints), and we added a lot of tests for real use
% cases during a long period of development and trial use.
% \end{response}

\item The backward/forward pass steps are not well described in the manuscript,
and figure 1b is hard to interpret. I suggest adding a supplemental figure
showing how these passes are performed. For example, the text explains that the
backward pass happens as information is passed to the root node -- does this
mean that the state is passed from `Mutagenize' to `Pool A' in or first from
`Pool Final' to `Stack'? And if so, what information is passed? Perhaps showing
each step in an expanded figure or labeling arrows with the order would be more
useful.

\begin{response}
The new Supplemental Information document includes a figure, Figure S1, that  breaks down the foward and backward pass computations in much more detail We believe this will answer the Reviewer's questions. 
\end{response}

% \begin{response}
% We agree that our description of state assignment (the backward pass) and
% sequence construction (the forward pass) was not clear. We have revised the
% Sequence generation subsection to specify the order of the two passes and what
% is passed between Pools and Operations. We have also expanded the example in
% Fig.~\ref{ms-fig:figure1}B into Additional file 1: Fig.~S1, numbering each step
% and showing the state
% calculations and sequence changes in action. We have revised the Sequence generation
% subsection as follows:
% \end{response}

% \begin{msquote}
% To generate a sequence library, PoolParty iterates through a series of
% \emph{states}. Each state is an integer that, together with a random seed,
% uniquely determines a sequence. The generation of each individual sequence
% proceeds in two steps: state assignment and sequence construction
% (Fig.~\ref{ms-fig:figure1}B; Additional file 1: Fig.~S1). This process begins when the corresponding
% state is passed to the root Pool. State assignment involves a
% \emph{backward pass} through the DAG, from the root Pool toward the source
% Pools. Each Pool passes the state it receives to the Operation that created it.
% Each Operation then decomposes this state into an \emph{internal state} (which
% it keeps) and zero or more \emph{component states}, each of which is passed to a
% distinct input Pool of that Operation. Sequence construction then occurs during
% a \emph{forward pass} through the DAG, from the initial Pools toward the root
% Pool: each Operation constructs a sequence based on its internal state and the
% sequences returned by its input Pools, then passes this sequence to the Pool it
% creates. The final constructed sequence is then received from the root Pool and
% recorded.
% \end{msquote}

\item A comparison to the pools produced by other tools such as VaLiAnT or
MPRAnator would provide additional confidence in the ability of PoolParty to
produce sequence pools as expected. Metrics such as the number of pool elements that overlap and an investigation of any unique pool elements would ensure that the generate pools are accurate.

\begin{response}
Figures S3 and S4 in the new Supplemental Information document compare the code/files needed to create identical libraries with PoolParty and VaLiAnT (Fig. S3) or tangermeme (Fig. S4). We attempted to make a similar figure for MPRNator, but that tool uses a web interface that is currently inoperative. 
\end{response}

\item Additional operations to improve library content would be useful. E.g. to maximize hamming distance, remove sequences containing restriction enzyme sites, remove homopolymers, etc. would make the tool more useful in generating ready-to-use pools.

\begin{response}
PoolParty does include a \texttt{filter} Operation that allows users to remove sequences that don't satisfy custom criterion. This Operation includes a number of built-in filter cirteria, including GC content, homopolymer content, sequence complexity, and restriction enzyme recodnition sites. We have also expanded the documentation on this capability and added a short description of this capability to the manuscript. 

\begin{msquote}
Sequences can be filtered based on sequence length, GC content, homopolymer
length, sequence complexity, restriction enzyme recognition sites, or other
user-defined criteria.
\end{msquote}

We have also clarified PoolParty's limitations in the Discussion section, in particular the fact that it not meant to generate libraries that are optimized for specific metrices.

\begin{msquote}
PoolParty is also not a sequence optimization tool; it does not design sequences that have optimized codon usage \citep{Swainston2017rb}, that are subject to complex synthesis constraints \citep{Zulkower2020jk}, or that have optimal activity as assessed by machine learning models \citep{Schreiber2025ledidi}.
\end{msquote}
\end{response}

% PoolParty does have a general \texttt{filter} Operation that lets users provide a
% Python function to decide
% which sequences should be kept. It also has several convenient filter methods
% for basic sequence properties, including GC content, homopolymers, sequence
% complexity, and restriction enzyme recognition sites. We realized that these
% features were not clearly described in the manuscript or documentation, and
% have now added a short description to the manuscript and expanded the
% documentation. We also clarified PoolParty's limitation in the Discussion to
% distinguish filtering from sequence optimization. Maximizing Hamming distance
% and other constraint-based optimizations usually require comparing sequences
% across the Pool and are outside the current scope of PoolParty. We have added the
% following sentence to the Pools and Operations subsection:
% \end{response}

% \begin{msquote}
% Sequences can be filtered based on sequence length, GC content, homopolymer
% length, sequence complexity, restriction enzyme recognition sites, or other
% user-defined criteria.
% \end{msquote}

% \begin{response}
% We have also revised the following sentence in the limitations paragraph of the
% Discussion:
% \end{response}

% \begin{msquote}
% PoolParty is also not a sequence optimization tool; it does not design sequences
% that have optimized codon usage \citep{Swainston2017rb}, that are subject to
% complex synthesis constraints \citep{Zulkower2020jk}, or that have optimal
% activity as assessed by machine learning models \citep{Schreiber2025ledidi}.
% \end{msquote}

\item The name PoolParty may be confused with the existing PoolParty analysis
softwre (\url{https://doi.org/10.1111/1755-0998.12784}).

\begin{response}
This is a fair concern, but we have decided to retain the name PoolParty. The package already uses the name \texttt{poolparty} on PyPI and ReadTheDocs, and announced this software in a preprint. Moreover, the previously published PoolParty is a pipeline for aligning and analyzing pooled-sequencing data, a use case that we believe is different enough from library sequence generation to avoid confusion with our package. 
\end{response}

% \begin{response}
% We thank the reviewer for pointing this out. The previously published PoolParty
% is a pipeline for aligning and analyzing pooled-sequencing data, whereas our
% package is used to design DNA sequence libraries. The two tools are unrelated
% and have different purposes. Given this difference in scope, we have retained
% the name PoolParty.
% \end{response}
\end{enumerate}

\reviewersection{Reviewer 3}

\begin{enumerate}
\item I think that the emphasis made on PoolParty being a universal tool that
can ``replace'' (line 256) existing code/tools needs to be qualified more. The
benefits of PoolParty are that it is straightforward to use, interpretable,
modular, with a nice `design card' tracking and graphed interpretation -- I
think this is particularly useful for ML and in silico work.

However, whist it does streamline design processes, the package operates on
input sequences rather than genomic loci co-ordinates/transcript so in terms of
wet-lab experiments it is most useful for fundamental biology cDNA DMS and MPRA
assays that are episomal (i.e. from plasmid / using reporters) rather than MAVEs
that rely upon genome editing (which is becoming the gold standard for clinical
variant work). It does not intuitively and natively support reference genome
coordinates, transcript models, exon/intron structures, alternative isoforms,
or genome assembly awareness (as far as I can tell); VaLiAnT does do this. Base
sequences into which variants are installed are represented/input as sequence
string (pasted by user) rather than standardized co-ordinate or HGVS notation
(this limits interpretation/ingress of variants from clinical and genomics
resources). If PoolParty is presented as intuitive and adaptable rather than a
universal replacement to alternative tools and code this is less of a concern
and may allow researchers to more readily select which software to invest time
into learning/using.

\begin{response}
This is a fair critique: PoolParty does not currently operate on genomes + transcript annotations, accept HGVS input, etc. We have removed language suggesting that PoolParty is meant to replace these other tools, including those that do, and have revised the Background and Discussion to define the scope of PoolParty more clearly. In Additional file 1, we have added two tables explicitly comparing the features and capabilities of PoolParty with existing tools (Tables~S2 and S3), and two figures showing side-by-side implementations of the same libraries in PoolParty versus VaLiAnT or tangermeme (Figs.~S3 and S4). Please also see the response to Comment 3 below.  
\end{response}

% \begin{msquote}
% Existing tools can help with parts of this process.
% \end{msquote}

% \begin{msquote}
% We have presented PoolParty, a Python package for designing oligonucleotide
% libraries. PoolParty facilitates this task by representing libraries as directed
% acyclic graphs (DAGs) that users can construct using a simple API.
% \end{msquote}

% \begin{msquote}
% PoolParty complements existing library-design tools such as VaLiAnT
% \citep{Barbon2022th} and MPRAnator \citep{GeorgakopoulosSoares2017gb} by
% providing a single framework for DMS libraries, MPRA libraries, and library
% designs that combine elements of both.
% \end{msquote}

% \begin{msquote}
% Additional comparisons between PoolParty and existing tools are provided in
% Additional file 1: Figs.~S3 and S4 and Tables~S2 and S3.
% \end{msquote}

\item Molecular consequence beyond codon change for missense is not reported as
far as I can tell, VaLiAnT does this to some extent (e.g. stop-gained,
synonymous, inframe deletion, etc, PoolParty may do this also, but I couldn't
see this after running a `replacement\_scan'?) but generally researchers use VEP
on their tool outputs to achieve extensive annotation. VEP requires standard
annotation or format (e.g. VCF) for input rather than sequence strings -- if
this is possible for PoolParty (the production of a VCF or a format that VEP
will accept) then one could obtain extensive annotations (e.g. SpliceAI
predictions, EVE, AlphaMissense, PolyPhen, Consequence, genomic coordinates,
Clinvar and gnomAD identifiers etc etc) which I think would greatly increase the
uptake of the tool by researchers in the MAVE community that are more
clinically/genomically minded -- if this is possible I would state this/provide
an easy means to generate VEP input format from PoolParty outputs.

\begin{response}
Yes, PoolParty does not annotate the molecular consequences of mutations, as it does not return an annotation layer with each sequence. However, if users wish to generate sequences that have a specific type of mutation, then fold these sequences into a larger library, there are two ways to track which specific sequences recieved which types of mutations: through the design card provided with each sequence, and through each sequence's name.

We appreciate the suggestion to support VCF format. To this end we have implemented a new method, \texttt{from\_vcf}, which extracts a specified sequence window and inserts the alternative alleles into that sequence. We do not, however, see a simple way of exporting libraries in VCF format, as the sequences generated by PoolParty are often not closely related to any specific wild-type sequence. The best option in this case would be for users to take the sequences generated by PoolParty, align them to a reference, and use the results to generate a correpsonding VCF record. We believe that this is best done using other available software, however, as building this capability into PoolParty would substantially expand PoolParty's dependencies.
\end{response}

% \begin{response}
% [REVISE THIS LATER]

% Thank you for raising these points. PoolParty doesn't directly annotate
% molecular consequences for variants, but for ORF-aware Operations these can
% usually be inferred from the Operation used and the amino-acid changes recorded
% in the design cards (i.e., \texttt{wt\_aas}, \texttt{mut\_aas}). In some cases,
% the consequence is already specified by the user, e.g., when requesting
% missense-only, synonymous, or nonsense mutations using
% \texttt{mutagenize\_orf}. We have also implemented ORF-aware insertion and
% deletion scans, with the insertion scan also supporting codon replacement,
% because the corresponding sequence-level Operations do not provide coding-aware
% design cards. These are local coding consequences only, though, since PoolParty
% does not track transcript context.

% The suggestion to support VCF output is great. We can imagine two ways of
% implementing it. The first is to make PoolParty aware of genomic coordinates,
% such that every Operation can track coordinate changes and propagate them
% downstream. The difficulty is that many Operations can't do this cleanly
% because the coordinate changes are not well defined, e.g., \texttt{join},
% \texttt{shuffle}, and \texttt{recombine}. The second is to implement an
% alignment algorithm that calls variants from the generated sequences post hoc,
% but this can be tricky for the same reason. It should be possible to support VCF
% output for a restricted set of Operations with well-defined coordinate changes,
% but we decided to leave this to a future version because it is not obvious what
% restrictions would be needed to prevent seemingly valid VCF output in
% unsupported cases. As an alternative, we implemented \texttt{from\_vcf}, which
% extracts the corresponding sequence window and inserts the alt allele
% independently for each VCF record. The resulting sequences can then be used as
% input to genomic deep learning models or for clinically focused episomal MAVE
% assays.
% \end{response}

\item The authors note limitations at the end of the discussion. I think
updating this section to address that PoolParty is mostly of benefit for DMS and
MPRA rather than genome editing MAVE assays (e.g. not as useful for Saturation
Genome Editing HDR libraries and Saturation Prime Editing pegRNA libraries)
which require extensive and specific annotation/standardized nomenclature and
consideration of transcript (generally through ingress of GTF/GFF files) and
exon/intron boundaries (including split codons between exons, again through
GTF/GFF boundary annotations).

Alternatively/additionally the authors might wish to highlight that the tool is
agnostic to species/genome/transcript and present it as a more intuitive
sequence generation tool which is particularly useful for ML
testing/benchmarking which is well tracked with the design cards and
interpretable/iteratively mutable through graph interrogation. This builds on
the comment that generally annotations are variant specific (line 56-58), with
the emphasis being that PoolParty serves as a useful tool to manifest oligo
libraries as structured objects (less of a concern for wet-lab investigations
in my opinion, but much more salient for ML based studies on raw sequence out of
genomic context).

\begin{response}
We agree and have updated the Limitations section accordingly. The revised textreads:
\end{response}

\begin{msquote} 
PoolParty operates on sequences, not genomic loci or transcript models. It can accept genomic coordinates as input to extract sequences, but does not interpret transcript annotations or keep track of those coordinates through subsequent design steps (e.g., as VaLiAnT does), which may make PoolParty less helpful for certain genome-editing applications.
\end{msquote}

\item Line 12 -- there is not a ``lack'' of purpose-built software, replace
with ``scarcity'' perhaps.

\begin{response}
Fair point. We have replaced ``lack'' with ``scarcity.'' The revised sentence reads:
\end{response}

\begin{msquote}
Designing these libraries, however, remains tedious and error-prone due to the scarcity of purpose-built software.
\end{msquote}

\item Line 51-53 - stated that VaLinAnT is assay specific, then say used for DMS
and saturation genome editing two different assays (SGE is a type of DMS, but
VaLinAnT is used for SGE and cDNA DMS which are quite different assays).

\begin{response}
We have replaced ``assay-specific'' with more general wording and revised the surrounding description. We also tested VaLiAnT on noncoding regions and confirmed that it supports these designs. We have revised the description as follows:
\end{response}

\begin{msquote}
Several other tools automate library design or sequence perturbation for particular applications. For example, VaLiAnT \citep{Barbon2022th} generates mutational scanning libraries of both coding and noncoding regions; supported mutation types include single-nucleotide substitutions, amino acid substitutions, and in-frame deletions.
\end{msquote}

\item Line 66 -- ``expensive computation'' this is only true for truly massive
sequence generation -- most full gene DMS/saturation libraries of $>$20,000
variants take $<$5 minutes on a standard laptop CPU (true of Valiant) -- this
statement would be true for many permutations of sequence for use in ML (e.g.
pairwise, full genome CDS on GPU) -- I would qualify this point on cost.

\begin{response}
Indeed, most libraries designed for wet-lab experiments can be generated in seconds to minutes on a standard CPU, as the reviewer notes. We have therefore made the wording more specific:
\end{response}

\begin{msquote}
Importantly, no sequences are generated until the user requests them, and it is simple to generate small test sets of sequences. This allows users to explore multiple design options without having to generate complete libraries at each iteration.
\end{msquote}

\item The use cases of protein coding, MPRA promoter and splice are good examples
(line 68).

\begin{response}
Thanks!
\end{response}

\item Figure 2 B and Figure 3 B -- is it necessary helpful to put these code
blocks here? I'm not sure how useful this is for readership.

\begin{response}
The purpose of these panels is to support our claim that the PoolParty API is simple to use. We believe that showing the code explicitly, together with the output the code generates, is the best way to support this claim, and so have opted to retain these panels. 
\end{response}

\item Line 194 -- `most abundant codon' -- this needs to be defined, is this
codon usage (I think not) or instance in the input sequence?

\begin{response}
We have clarified the text:
\end{response}

\begin{msquote}
We note that \texttt{mutagenize\_orf} supports multiple codon mutation types; this example uses \texttt{missense\_only\_first}, which selects the most frequently used codon in the human genome for each target amino acid.
\end{msquote}

\item Line 235 -- SpliceAI is defined here for the first time but is introduced
before this in section 2.8 -- I would move the `a deep learning model of mRNA
splicing' to the earlier section.

\begin{response}
The term SpliceAI is now defined in the earlier Implementation subsection, where the
this term is first used:
\end{response}

\begin{msquote}
SpliceAI \citep{Jaganathan2019ps} is a genomic deep learning model for splice-site prediction in humans. We used the five-model ensemble to predict the canonical 5'ss probability using $\pm$5~kb of genomic context (GRCh38).
\end{msquote}

\item Line 238 ``varying strength'' -- maxentscan in legend, I would define this
is the main body text too.

\begin{response}
We have clarified that cryptic splice-site strength refers to the MaxEntScan. The revised sentence in Results reads:
\end{response}

\begin{msquote}
Using the 5'ss of SMN2 exon~7 as the canonical site, we used PoolParty to construct a library in which cryptic 5'ss 9-mers of varying strength
(MaxEntScan score) were inserted at varying positions on either side of the canonical 5'ss (see Implementation).
\end{msquote}

\item Line 248 -- Fig.4G is introduced before 4E-F, perhaps generally introduce
Fig.4. then this could stay as is.

\begin{response}
We have moved the cross-validation result to the end of the paragraph so that
the Figure 4 panels are introduced in order. 
\end{response}

\item Line 256 -- ``replaces the ad hoc'' is too bold and a little presumptive as
for major points above-- needs clarification/qualification. Also, Line 274 --
``in contrast''-- I think these claims are too strong -- valiant can do DMS and
MPRA as well as SGE/SPE -- I would tone these statements down and/or emphasise
what is unique to PoolParty.

\begin{response}
This is fair criticism; we have revised both statements:
\end{response}

\begin{msquote}
We have presented PoolParty, a Python package for designing oligonucleotide libraries. PoolParty facilitates this task by representing libraries as directed acyclic graphs (DAGs) that users can construct using a simple API.
\end{msquote}

\begin{msquote}
PoolParty complements existing library-design tools such as VaLiAnT \citep{Barbon2022th} and MPRAnator \citep{GeorgakopoulosSoares2017gb} by
providing a single framework for DMS libraries, MPRA libraries, and library designs that combine elements of both.
\end{msquote}

\item Line 271 -- ``in the spliceAI example'' -- the one in section 2.8 or 3.3?
(3.3 I imagine, but best to state).

\begin{response}
We have clarified that this refers to the SpliceAI surrogate modeling example in the Results:
\end{response}

\begin{msquote}
For example, in the SpliceAI analysis of Fig.~\ref{ms-fig:figure4}, the cryptic splice-site strength and position recorded in each design card provided the covariates used for surrogate modeling.
\end{msquote}

\item The code is very well written and documented, and it is noted by the
authors that Anthropic models have been used to do this/help which I think is
sensible, editorially this may be something to confirm is acceptable at
Bioinformatics (as reviewers were requested not to use these tools, but this is
likely due to confidentiality).

\begin{response}
BMC Bioinformatics policy permits LLM use when disclosed. We have also clarified the LLM disclosure statement in Implementation:
\begin{msquote}
The authors used large language models (primarily Claude Opus 4.5 and 4.6, Anthropic) to assist with code development, debugging, and documentation writing. The manuscript was written entirely by the authors, with LLM assistance only at the copy-editing stage. The figures were also prepared by the authors; LLMs were used only to help write the scripts that generate the plots. The authors affirm that they are fully responsible for the content of the manuscript, the codebase, and the documentation. 
\end{msquote}
\end{response}

\item The name `PoolParty' is great.

\begin{response}
Thank you!
\end{response}
\end{enumerate}

\bibliographystyle{sn-vancouver-num}
\bibliography{references}

\end{document}
